\subsection{形式光滑局部代数的一个判别法}

命题 7. 设 \(k_{0}\) 为环，\(k\) 为 \(k_{0}\)-代数，A 为 \(k\)-代数，\(\mathfrak{m}\) 为 A 的一个极大理想。假设 \(k\) 与 A/\(\mathfrak{m}\) 都在 \(k_{0}\) 上形式光滑。为使 A 在 \(k\) 上对 \(\mathfrak{m}\)-进拓扑形式光滑，必须且只需下列两个条件成立：

\begin{enumerate}
\def\labelenumi{(\roman{enumi})}
\item
  典范同态 \(\mathsf{S}_{\mathrm{A} / \mathfrak{m}}(\mathfrak{m} / \mathfrak{m}^{2})\to \mathrm{gr}_{\mathfrak{m}}(\mathrm{A})\) 是双射；
\item
  A/m-线性映射
\end{enumerate}

\[
\omega : \mathrm{A} / \mathfrak {m} \otimes_ {k} \Omega_ {k _ {0}} (k) \longrightarrow \mathrm{A} / \mathfrak {m} \otimes_ {\mathrm{A}} \Omega_ {k _ {0}} (\mathrm{A})
\]

由典范映射 \(k \to A\) 导出，是单射。

设 \(d_{k}:k\to \Omega_{k_{0}}(k)\) 与 \(d_{\mathrm{A}}:\mathrm{A}\to \Omega_{k_0}(\mathrm{A})\) 为泛 \(k_{0}\)-导子。

先设 A 在 \(k\) 上对 \(\mathfrak{m}\)-进拓扑形式光滑。则 A 在 \(k_{0}\) 上对 \(\mathfrak{m}\)-进拓扑形式光滑（\(n^{\circ} 2\), 命题 3, a)），这等价于 (i)（\(n^{\circ} 5\), 定理 2）。此外，\(k_{0}\)-导子 \(\lambda \mapsto 1 \otimes d_{k}(\lambda)\)（\(k\) 的、取值于 \(A/\mathfrak{m} \otimes_{k} \Omega_{k_{0}}(k)\) 中的）可以扩张为一个 \(k_{0}\)-导子（A 的、取值于 \(A/\mathfrak{m} \otimes_{k} \Omega_{k_{0}}(k)\) 中的）（\(n^{\circ} 2\), 注 2）。于是存在 A-线性映射 \(u: \Omega_{k_{0}}(A) \to A/\mathfrak{m} \otimes_{k} \Omega_{k_{0}}(k)\)，满足 \(u(d_{A}(\lambda 1_{A})) = 1 \otimes d_{k}(\lambda)\)（对每个 \(\lambda \in k\)）。\(A/\mathfrak{m}\)-线性映射 \(A/\mathfrak{m} \otimes_{A} \Omega_{k_{0}}(A) \longrightarrow A/\mathfrak{m} \otimes_{k} \Omega_{k_{0}}(k)\)（由 \(u\) 导出）是 \(\omega\) 的一个收缩，这就证明了 (ii)。

反之，设条件 (i) 与 (ii) 成立。则 A 在 \(k_{0}\) 上对 \(\mathfrak{m}\)-进拓扑形式光滑（\(n^{\circ}5\), 定理 2），且 A-模 \(\Omega_{k_{0}}(A)\) 是投射的（\(n^{\circ}5\), 定理 2 的推论 3）。取定一个整数 \(r \geqslant 0\)，并考虑 \(A/\mathfrak{m}^{r}\)-线性映射

\[
\omega_ {r}: \mathrm{A} / \mathfrak {m} ^ {r} \otimes_ {k} \Omega_ {k _ {0}} (k) \longrightarrow \mathrm{A} / \mathfrak {m} ^ {r} \otimes_ {\mathrm{A}} \Omega_ {k _ {0}} (\mathrm{A})
\]

由典范映射 \(k\to\mathrm{A}\) 导出。设 \((\lambda_{i})_{i\in\mathrm{I}}\) 为 \(k\) 的一个元素族，使诸 \(d_{k}(\lambda_{i})\) 构成 \(k\)-向量空间 \(\Omega_{k_{0}}(k)\) 的一个基。由 (ii)，元素 \(1 \otimes d_{A}(\lambda_{i}1_{A})\) 在 A/m \(\otimes_{\mathrm{A}} \Omega_{k_0}(\mathrm{A})\) 中线性无关。据 II, § 3, 第 2 节, 命题 5 的推论 1 与 2，诸 \(1 \otimes d_{A}(\lambda_{i}1_{A})\) 构成 A/m \(^{r}\)-模 A/m \(^{r} \otimes_{\mathrm{A}} \Omega_{k_0}(\mathrm{A})\) 的一个直和项的基。于是存在一个 A/m \(^{r}\)-线性映射

\[
u _ {r}: \mathrm{A} / \mathfrak {m} ^ {r} \otimes_ {\mathrm{A}} \Omega_ {k _ {0}} (\mathrm{A}) \longrightarrow \mathrm{A} / \mathfrak {m} ^ {r} \otimes_ {k} \Omega_ {k _ {0}} (k)
\]

满足 \(u_r(1 \otimes d_A(\lambda_i 1_A)) = 1 \otimes d_k(\lambda_i)\)（对每个 \(i\)），因而 \(u_r \circ \omega_r = \mathrm{Id}\)。

现在来验证 A 在 \(k\) 上对 m-进拓扑形式光滑。设 C 为 \(k\)-代数，N 为 C 中的平方零理想，\(\pi: C \to C/N\) 为典范满射；赋予 C 与 C/N 以离散拓扑。设 \(\varphi: A \to C/N\) 为连续 \(k\)-代数同态。由于 A 在 \(k_0\) 上对 m-进拓扑形式光滑，存在连续 \(k_0\)-代数同态 \(\tilde{\varphi}_0: A \to C\)，满足 \(\pi \circ \tilde{\varphi}_0 = \varphi\)。据第 1 节的命题 1，\(k_0\)-代数同态 \(\tilde{\varphi}: A \to C\)（满足 \(\pi \circ \tilde{\varphi} = \varphi\) 者）恰是形如 \(x \mapsto v(d_A(x)) + \tilde{\varphi}_0(x)\) 的映射，其中 \(v\) 取遍 \(\mathrm{Hom}_{A}(\Omega_{k_0}(A), N)\)。尚需选取 \(v\)，使 \(\tilde{\varphi}\) 是 \(k\)-代数同态。映射 \(\lambda \mapsto \lambda 1_C - \tilde{\varphi}_0(\lambda 1_A)\) 是一个 \(k_0\)-导子（自 \(k\) 到 N 中）(同上引文)，从而可写成 \(h \circ d_k\)，其中 \(h \in \mathrm{Hom}_k(\Omega_{k_0}(k), N)\)。

取一个整数 \(r\) 使 \(\varphi\) 的核包含 \(\mathfrak{m}^r\)。A-模 N 被 \(\mathfrak{m}^r\) 零化，只需取 \(v\) 为下列同态序列的复合

\[
\Omega_ {k _ {0}} (\mathrm{A}) \longrightarrow \mathrm{A} / \mathfrak {m} ^ {r} \otimes_ {\mathrm{A}} \Omega_ {k _ {0}} (\mathrm{A}) \xrightarrow {u _ {r}} \mathrm{A} / \mathfrak {m} ^ {r} \otimes_ {k} \Omega_ {k _ {0}} (k) \xrightarrow {h ^ {\prime}} \mathrm{N},
\]

其中 \(h'\) 由 \(h\) 导出。实际上，我们有（对 \(\lambda \in k\)）：

\[
v \left(d _ {\mathrm{A}} \left(\lambda 1 _ {\mathrm{A}}\right)\right) = h ^ {\prime} u _ {r} \left(1 \otimes d _ {\mathrm{A}} \left(\lambda 1 _ {\mathrm{A}}\right)\right) = h ^ {\prime} \left(1 \otimes d _ {k} (\lambda)\right) = h \left(d _ {k} (\lambda)\right) = \lambda 1 _ {\mathrm{C}} - \tilde {\varphi} _ {0} \left(\lambda 1 _ {\mathrm{A}}\right).
\]

注 1.—当 A 为 Noether 环时，条件 (i) 表示局部环 \(A_{\mathfrak{m}}\) 是正则的 (VIII, § 5, 第 2 节, 定理 1)。

命题 8. 设 \(k\) 为域，A 为 Noether 局部 \(k\)-代数。下列条件等价：

\begin{enumerate}
\def\labelenumi{(\roman{enumi})}
\item
  A 在 \(k\) 上对 \(\mathfrak{m}_{A}\)-进拓扑形式光滑；
\item
  A 正则，且 \(\kappa_{\mathrm{A}}\)-线性映射
\end{enumerate}

\[
\omega : \kappa_ {\mathrm{A}} \otimes_ {k} \Omega (k) \longrightarrow \kappa_ {\mathrm{A}} \otimes_ {\mathrm{A}} \Omega (\mathrm{A})
\]

由典范单射 k → A 导出，是单射；

\begin{enumerate}
\def\labelenumi{(\roman{enumi})}
\setcounter{enumi}{2}
\item
  A 绝对正则；
\item
  对每个纯不可分扩张 \(k'\)（\(k\) 的、有限次且高度 \(\leqslant 1\) 的），局部环 \(A_{(k')}\) 是正则的。
\item
  (ii) \(\Leftrightarrow\) (i)：只需应用命题 7 与上面的注 1，取 \(k_{0}\) 为 \(k\) 的素子域；实际上，\(k\) 与 \(\kappa_{\mathrm{A}}\) 都在 \(k_{0}\) 上形式光滑（\(n^{\circ}3\), 参见定理 1）。
\end{enumerate}

(i)⇒(iii)：这由定理 2 的推论 2（第 5 节）得出。

(iii)⇒(iv)：这由 § 6 第 4 节的定义 1 得出。

若 \(k\) 的特征为 0，则由定理 2 的推论 1（第 5 节）可知 (iv) 蕴含 (i)，从而命题在此情形得证。设 \(k\) 的特征为 \(p \neq 0\)，证明 (iv)⇒(ii)。设 \(k'\) 为 \(k\) 的一个有限次且高度 \(\leqslant 1\) 的纯不可分扩张。若 A 与 \(\mathrm{A}_{(k')}\) 都正则，则典范映射 \(\kappa_{\mathrm{A}} \otimes_{k'^{p}} \Omega_{k^{p}}(k'^{p}) \longrightarrow \kappa_{\mathrm{A}} \otimes_{\mathrm{A}} \Omega(\mathrm{A})\) 是单射（第 6 节, 命题 6）。据 A, V, 第 97 页, 定理 1 b)（应用于扩张 \(k\)（\(k^{p}\) 上的）），\(k\)-向量空间 \(\Omega(k)\)（它与 \(\Omega_{k^{p}}(k)\) 重合）是诸子空间 \(k \otimes_{k'^{p}} \Omega_{k^{p}}(k'^{p})\) 的递增滤过并，其中 \(k'\) 取遍 \(k\) 的、高度 \(\leqslant 1\) 的有限纯不可分扩张（含于 \(k\) 的一个固定代数闭包中的）。断言 (ii) 由此得出。

注 2.—设 \(k\) 为域，A 为一个 \(k\)-代数，使得环 \(\mathrm{A}_{(k')}\) 对每个纯不可分扩张 \(k'\)（\(k\) 的、有限次且高度 \(\leqslant 1\) 的）都是正则的；则 A 绝对正则。实际上，设 \(k'\) 是这样一个扩张；对 A 的每个极大理想 \(\mathfrak{m}\)，环 \(k'\otimes_k\mathrm{A}_{\mathfrak{m}}\) 与 \(\mathrm{A}_{(k')}\) 的一个分式环恒同，因而是正则的。由上面的命题 8 与 § 6 第 4 节的命题 6，代数 A 绝对正则。

